\documentclass[a4,11pt]{aleph-notas}

% -- Paquetes adicionales
\usepackage{aleph-comandos}
\usepackage{booktabs}
\usepackage{pdflscape}
\hypersetup{
    urlcolor=colordef!70,
}

% -- Datos  
\institucion{Facultad de Ciencias Exactas, Naturales y Ambientales}
\carrera{Catálogo STEM}
\asignatura{Matemática III}
\tema{Reto no. 2: Análisis de Componentes Principales}
\autor{Andrés Merino}
\fecha{Periodo 2026-1}

\logouno[0.16\textwidth]{Logos/logoPUCE_04_ac}
\definecolor{colortext}{HTML}{1957d1}
\definecolor{colordef}{HTML}{1957d1}
\fuente{montserrat}


% -- Otros comandos




\begin{document}

\encabezado

%%%%%%%%%%%%%%%%%%%%%%%%%%%%%%%%%%%%%%%%
\section{Resultado de aprendizaje}
%%%%%%%%%%%%%%%%%%%%%%%%%%%%%%%%%%%%%%%%

\begin{itemize}
\item 
    Aplica fundamentos de Álgebra Lineal y técnicas de análisis multivariado para reducir dimensionalidad, interpretar estructuras de datos y comunicar hallazgos reproducibles en contextos de ciencia de datos.
\end{itemize}


%%%%%%%%%%%%%%%%%%%%%%%%%%%%%%%%%%%%%%%%
\section{Descripción}
%%%%%%%%%%%%%%%%%%%%%%%%%%%%%%%%%%%%%%%%

El análisis de componentes principales (ACP) es una técnica central en ciencia de datos para reducir dimensionalidad, eliminar redundancia y facilitar la interpretación de información multivariada. Este reto busca articular el fundamento matemático del ACP con su implementación computacional y su lectura analítica en un flujo de trabajo realista de análisis de datos.

\subsection*{Pregunta esencial}
\begin{itemize}[leftmargin=*]
\item ¿Cómo transforma el ACP un conjunto de variables correlacionadas en nuevas variables informativas para apoyar decisiones en ciencia de datos?
\end{itemize}

\subsection*{Reto}

Tu desafío es doble: primero, explicar con claridad la base matemática del ACP; segundo, aplicarlo a un conjunto de datos real para extraer conclusiones interpretables.

Debes entregar:
\begin{itemize}[leftmargin=*]
\item Un minitutorial en video (máx. 8 minutos) con explicación conceptual y procedimiento.
\item Un informe en PDF redactado en \LaTeX\ con análisis teórico y resultados prácticos.
\item Un notebook de Jupyter o Google Colab reproducible con el proceso completo.
\end{itemize}

\subsection*{Preguntas guía}
\begin{itemize}[leftmargin=*]
\item ¿Por qué el centrado y la estandarización de variables son pasos críticos antes del ACP?
\item ¿Cómo se interpreta la varianza explicada por cada componente principal?
\item ¿Qué relación existe entre autovalores/autovectores y la proyección ortogonal en el ACP?
\item ¿Cómo decidir entre usar 2 o 3 componentes según el objetivo analítico?
\item ¿Qué información aportan las cargas (loadings) para interpretar variables relevantes?
\end{itemize}

\subsection*{Actividades guía}
\begin{enumerate}[leftmargin=*, label={{\arabic*.}}]
\item Investigar y explicar ACP desde Álgebra Lineal: base, dimensión, transformación lineal, ortogonalidad, autovalores y autovectores.
\item Elaborar un guion para un video breve, claro y didáctico.
\item Descargar el dataset \href{https://www.kaggle.com/datasets/konradb/real-world-vehicle-emissions}{Real-world vehicle emissions} y usar el archivo \texttt{2021 Cars Raw.csv}.
\item Realizar preprocesamiento mínimo: selección de variables cuantitativas, limpieza de faltantes y estandarización.
\item Aplicar ACP con \texttt{scikit-learn} para 2 y 3 componentes.
\item Presentar visualizaciones: dispersión en espacio de componentes, varianza explicada y análisis de cargas.
\item Verificar manualmente la proyección de al menos un registro usando la base ortonormal asociada a los componentes seleccionados.
\item Redactar conclusiones enfocadas en hallazgos, utilidad del método y limitaciones del análisis.
\end{enumerate}

\subsection*{Recursos}
\begin{itemize}[leftmargin=*]
\item Video \href{https://youtu.be/FgakZw6K1QQ?si=a9cxs4ePio-H6MXm}{StatQuest: Análisis de componentes principales (PCA), paso a paso}.
\item Video \href{https://youtu.be/7My_PBhxeP4?si=zQf5Hr18NOb6cerX}{Análisis de componentes principales (PCA)}.
\item Documentación oficial de \href{https://scikit-learn.org/stable/modules/generated/sklearn.decomposition.PCA.html}{\texttt{sklearn.decomposition.PCA}}.
\item Complementar con al menos dos fuentes académicas (libro o artículo) citadas en formato APA.
\end{itemize}

%%%%%%%%%%%%%%%%%%%%%%%%%%%%%%%%%%%%%%%%
\section{Producto}
%%%%%%%%%%%%%%%%%%%%%%%%%%%%%%%%%%%%%%%%

\subsection{Minitutorial explicativo}
Un video de máximo 8 minutos (formato libre) donde se explique con claridad:
\begin{itemize}[leftmargin=*]
\item Qué es el ACP y para qué se usa en ciencia de datos.
\item Relación del ACP con base, dimensión, proyección ortogonal y autovectores.
\item Flujo mínimo de implementación: preprocesamiento, ajuste, transformación e interpretación.
\item Ejemplo breve de lectura de resultados (varianza explicada y gráfico de componentes).
\end{itemize}

\subsection{Informe de aplicación práctica}
Un documento en formato PDF que contenga:

\begin{enumerate}[leftmargin=*, label={\textbf{\arabic*.}}]
\item \textbf{Introducción:} contexto del problema, objetivo analítico y pertinencia del ACP.
\item \textbf{Fundamento matemático:} explicación conceptual de los elementos de Álgebra Lineal implicados.
\item \textbf{Desarrollo computacional:} limpieza básica, selección de variables y ejecución de ACP (2 y 3 componentes) con \texttt{scikit-learn}.
\item \textbf{Resultados e interpretación:} visualizaciones, varianza explicada, cargas y análisis de hallazgos.
\item \textbf{Verificación manual:} proyección de al menos un dato sobre la base de componentes seleccionada.
\item \textbf{Cierre y bibliografía:} conclusiones, limitaciones y referencias en formato APA.
\end{enumerate}

\subsection{Entrega técnica}
\begin{itemize}[leftmargin=*]
\item El informe debe ser elaborado en la plantilla \LaTeX\ institucional disponible en Overleaf: \href{https://www.overleaf.com/latex/templates/formato-tareas-puce/nkgwqjtcrvms}{Formato Tareas PUCE}.
\item Incluir el enlace al video con permisos de acceso activos.
\item Incluir el enlace al notebook (Google Colab o repositorio) con celdas ejecutables.
\item Si los enlaces no permiten visualizar el material, la actividad no podrá ser calificada y se asignará \textbf{0}.
\end{itemize}

%%%%%%%%%%%%%%%%%%%%%%%%%%%%%%%%%%%%%%%%
\section{Rúbrica de calificación}
%%%%%%%%%%%%%%%%%%%%%%%%%%%%%%%%%%%%%%%%

En caso de presentar un video con voz en off, rostro no visible o fuera del rango de tiempo establecido, se reducirá un nivel a cada punto de la rúbrica.

\begin{landscape}
\begin{center}\footnotesize
\begin{tabular}{p{5cm}p{4cm}p{4cm}p{4cm}p{4cm}}
\toprule
\textbf{Indicador} & \textbf{Excelente} & \textbf{Bueno} & \textbf{Aceptable} & \textbf{Insuficiente} \\ \midrule
\textbf{Fundamento matemático del ACP (10 pts)} & Explica con rigor base, dimensión, proyección ortogonal y autovectores vinculados al ACP, con notación correcta (10) & Explicación correcta con una omisión menor conceptual o de notación (8) & Presenta ideas generales, pero con justificación parcial o confusa (6) & Presenta errores conceptuales importantes o no justifica (0) \\ \midrule
\textbf{Implementación del ACP en notebook (15 pts)} & Preprocesa adecuadamente, ejecuta ACP para 2 y 3 componentes, y documenta el flujo de forma reproducible (15) & Implementación correcta con detalles menores faltantes (12) & Implementación parcial o con pasos incompletos (9) & Errores técnicos relevantes o sin evidencia funcional (0) \\ \midrule
\textbf{Interpretación de resultados (15 pts)} & Interpreta correctamente varianza explicada, cargas y gráficos; incluye verificación manual de proyección (15) & Interpretación correcta con omisiones menores en el análisis (12) & Interpretación parcial o poco conectada con el problema (9) & Interpretación incorrecta o ausente (0) \\ \midrule
\textbf{Presentación en video (5 pts)}
& Exposición clara y fluida; rostro visible durante toda la presentación (5)
& Exposición clara con ligeras interrupciones; rostro visible (4)
& Exposición poco fluida o lectura excesiva; rostro visible parcial (3)
& Voz en off o rostro no visible / exposición ininteligible (0) \\ \midrule
\textbf{Gestión del tiempo (5 pts)} & 8:00 $\pm$ 30 s (5) & 8:00 $\pm$ 60 s (4) & 8:00 $\pm$ 90 s (3) & Fuera de $\pm$ 90 s (0) \\ \bottomrule
\end{tabular}
\end{center}
\end{landscape}

\end{document}